
\documentclass{article}
\usepackage{colortbl}
\usepackage{makecell}
\usepackage{multirow}
\usepackage{supertabular}

\begin{document}

\newcounter{utterance}

\centering \large Interaction Transcript for game `cladder', experiment `full\_v1.5\_default', episode 3258 with llama{-}3.1{-}8b{-}instruct.
\vspace{24pt}

{ \footnotesize  \setcounter{utterance}{1}
\setlength{\tabcolsep}{0pt}
\begin{supertabular}{c@{$\;$}|p{.15\linewidth}@{}p{.15\linewidth}p{.15\linewidth}p{.15\linewidth}p{.15\linewidth}p{.15\linewidth}}
    \# & \multicolumn{2}{c}{Player} && \multicolumn{2}{c}{Game Master} \\
    \hline

    \theutterance \stepcounter{utterance}  
    & & & \multicolumn{4}{p{0.6\linewidth}}{
        \cellcolor[rgb]{0.9,0.9,0.9}{
            \makecell[{{p{\linewidth}}}]{
                \texttt{\tiny{[P1$\langle$GM]}}
                \texttt{You are an expert in causal inference. The following question is not a typical commonsense query, but rather a meticulously designed question created by a professor specializing in causal inference, intended to assess the students' mastery of the course content.} \\
\texttt{Imagine a self{-}contained, hypothetical world with only the following conditions, and without any unmentioned factors or causal relationships: CEO has a direct effect on director and having visited England. Having visited England has a direct effect on employee. Director has a direct effect on employee. The overall probability of having visited England is 19\%. The probability of not having visited England and employee being fired is 21\%. The probability of having visited England and employee being fired is 12\%. Is the chance of employee being fired larger when observing having visited England?} \\
\texttt{Start your answer with 'Yes' or 'No', followed by additional reasoning or evidence to support your explanation.} \\
            }
        }
    }
     \\ \\

    \theutterance \stepcounter{utterance}  
    & \multicolumn{4}{p{0.6\linewidth}}{
        \cellcolor[rgb]{0.561,0.737,0.561}{
            \makecell[{{p{\linewidth}}}]{
                \texttt{\tiny{[P1$\rangle$GM]}}
                \texttt{No.} \\
\\ 
\texttt{To answer this question, we need to apply the principles of causal inference, specifically the concept of confounding variables and the potential outcome framework. } \\
\\ 
\texttt{Let's denote the variables as follows: } \\
\\ 
\texttt{{-} E: Employee being fired} \\
\texttt{{-} V: Having visited England} \\
\texttt{{-} D: Director} \\
\texttt{{-} C: CEO} \\
\\ 
\texttt{We are given the following information:} \\
\\ 
\texttt{1. P(V) = 0.19 (probability of having visited England)} \\
\texttt{2. P(E|\textasciitilde{}V, \textasciitilde{}D) = 0.21 (probability of employee being fired given not having visited England and not being a director)} \\
\texttt{3. P(E|V, \textasciitilde{}D) = 0.12 (probability of employee being fired given having visited England and not being a director)} \\
\\ 
\texttt{We are also given the following causal relationships:} \\
\\ 
\texttt{{-} C → D} \\
\texttt{{-} C → V} \\
\texttt{{-} V → E} \\
\texttt{{-} D → E} \\
\\ 
\texttt{To determine whether the chance of employee being fired is larger when observing having visited England, we need to compare the probability of employee being fired given having visited England (P(E|V)) to the probability of employee being fired given not having visited England (P(E|\textasciitilde{}V)).} \\
\\ 
\texttt{However, we are not given the direct probability of employee being fired given having visited England (P(E|V)) and being a director (P(D|V)). We only have the probability of employee being fired given having visited England and not being a director (P(E|V, \textasciitilde{}D)).} \\
\\ 
\texttt{To address this, we can use the law of total probability to calculate the probability of employee being fired given having visited England (P(E|V)):} \\
\\ 
\texttt{P(E|V) = P(E|V, D)P(D|V) + P(E|V, \textasciitilde{}D)P(\textasciitilde{}D|V)} \\
\\ 
\texttt{We are not given the probability of being a director given having visited England (P(D|V)), but we can use the given information to derive it.} \\
\\ 
\texttt{Since the probability of having visited England is 0.19, the probability of not having visited England is 0.81. We can use the law of total probability to calculate the probability of being a director given having visited England (P(D|V)):} \\
\\ 
\texttt{P(D|V) = P(D|V, E)P(E|V) + P(D|V, \textasciitilde{}E)P(\textasciitilde{}E|V)} \\
\\ 
\texttt{However, we are not given the probability of employee being fired given having visited England and being a director (P(E|V, D)) or the probability of employee not being fired given having visited England and being a director (P(\textasciitilde{}E|V, D)). We only have the probability of employee being fired given having visited England and not being a director (P(E|V, \textasciitilde{}D)).} \\
\\ 
\texttt{To address this, we can use the given information to derive the probability of employee being fired given having visited England and being a director (P(E|V, D)).} \\
\\ 
\texttt{We are given the} \\
            }
        }
    }
    & & \\ \\

    \theutterance \stepcounter{utterance}  
    & & & \multicolumn{2}{p{0.3\linewidth}}{
        \cellcolor[rgb]{0.9,0.9,0.9}{
            \makecell[{{p{\linewidth}}}]{
                \texttt{\tiny{[GM$|$GM]}}
                \texttt{No} \\
            }
        }
    }
    & & \\ \\

    \theutterance \stepcounter{utterance}  
    & & & \multicolumn{2}{p{0.3\linewidth}}{
        \cellcolor[rgb]{0.9,0.9,0.9}{
            \makecell[{{p{\linewidth}}}]{
                \texttt{\tiny{[GM$|$GM]}}
                \texttt{yes} \\
            }
        }
    }
    & & \\ \\

    \theutterance \stepcounter{utterance}  
    & & & \multicolumn{2}{p{0.3\linewidth}}{
        \cellcolor[rgb]{0.9,0.9,0.9}{
            \makecell[{{p{\linewidth}}}]{
                \texttt{\tiny{[GM$|$GM]}}
                \texttt{game\_result = LOSE} \\
            }
        }
    }
    & & \\ \\

\end{supertabular}
}

\end{document}
